\documentclass[12pt,a4paper]{article}

% ===== 基本包 =====
\usepackage[utf8]{inputenc}
\usepackage[T1]{fontenc}
\usepackage{lmodern}
\usepackage{amsmath,amssymb,amsfonts}
\usepackage{graphicx}
\usepackage{booktabs}
\usepackage{caption}
\usepackage{subcaption}

% ===== 格式设置 =====
\usepackage[margin=2.5cm]{geometry}
\usepackage{setspace}
\onehalfspacing
\usepackage[parfill]{parskip}

% ===== 参考文献管理 =====
\usepackage[backend=biber,style=authoryear-comp,natbib=true]{biblatex}
\addbibresource{../../references/library.bib}

% ===== 标题信息 =====
\title{[论文标题]}
\author{作者姓名\\[2mm]
  \small 单位名称 \\
  \small \texttt{email@example.com}
}
\date{\today}

% ===== 自定义命令 =====
\newcommand{\sectionref}[1]{Section~\ref{#1}}
\newcommand{\equationref}[1]{Equation~(\ref{#1})}
\newcommand{\figureref}[1]{Figure~\ref{#1}}
\newcommand{\tableref}[1]{Table~\ref{#1}}

\begin{document}

\maketitle

\begin{abstract}
  [请在此处撰写摘要，简要说明研究背景、方法、主要发现和贡献。]

  \vspace{2mm}
  \textbf{关键词：} [关键词1], [关键词2], [关键词3]
\end{abstract}

% ===== 1. 引言 =====
\section{引言}
\label{sec:introduction}

研究背景和动机...
主要贡献概述...

% ===== 2. 相关工作 =====
\section{相关工作}
\label{sec:related_work}

文献综述...
现有方法的局限性...

% ===== 3. 方法 =====
\section{提出的方法}
\label{sec:methodology}

\subsection{方法概述}
\label{subsec:overview}

\subsection{具体实现}
\label{subsec:implementation}

\subsection{算法复杂度分析}
\label{subsec:complexity}

% ===== 4. 实验结果 =====
\section{实验}
\label{sec:experiments}

\subsection{实验设置}
\label{subsec:setup}

数据集描述...
评估指标...

\subsection{结果与分析}
\label{subsec:results}

\begin{figure}[htbp]
  \centering
  \includegraphics[width=0.8\textwidth]{../../figures/example.pdf}
  \caption{实验结果示意图}
  \label{fig:results}
\end{figure}

\begin{table}[htbp]
  \centering
  \caption{性能对比}
  \label{tab:comparison}
  \begin{tabular}{lccc}
    \toprule
    方法 & 精度 & 召回率 & F1分数 \\
    \midrule
    方法1 & 0.85 & 0.82 & 0.83 \\
    方法2 & 0.87 & 0.84 & 0.85 \\
    我们的方法 & \textbf{0.91} & \textbf{0.89} & \textbf{0.90} \\
    \bottomrule
  \end{tabular}
\end{table}

% ===== 5. 讨论 =====
\section{讨论}
\label{sec:discussion}

结果的深入分析...
方法的优势与局限性...

% ===== 6. 结论 =====
\section{结论}
\label{sec:conclusion}

总结主要发现...
展望未来工作...

% ===== 参考文献 =====
\printbibliography

% ===== 附录 =====
\appendix
\section{附录}
\label{sec:appendix}

补充材料...

\end{document}